\specialchapter{引言}

谱理论旨在研究某些赋范代数的性质，其基本的例子是：拓扑空间上复值连续函数所成的代数、Banach 空间的自同态所成的代数，以及关于局部紧群上 Haar 测度可积的函数所成的卷积代数。仅由这一简单的说明未必能看出这项研究的范围，但要认识其重要性，只需注意：两位当代数学家若不属于同一学派，则除了线性代数以及微分学与积分学的初等部分之外，他们共同具有的知识往往只是 Fourier 分析的某些方面。

在本书的头两章中，我们展开谱理论中最基本的思想，并更详细地考察上面提到的第三个例子，即局部紧交换群框架下的 Fourier 变换。

随后的各章将首先继续紧线性映射的一般理论，然后是 Hilbert 空间上算子的理论，后一情形还包括仅定义在这类空间的一个子空间（通常是稠密子空间）上的线性映射。众所周知，这一理论是量子力学数学形式体系的基础。最后，我们将给出拓扑群的酉表示的初等理论，特别是紧群的酉表示；在《李群与李代数》一书第九章中我们已有机会用到它。

头两章的标题与本书 1967 年出版的第一版相同。它们经过了修订、更正与扩充。

为便于阅读，我们在下面指出本版中最重要的改动，特别是增补的内容。

关于第 I 章：

- 在 §3 的第 3 小节中，我们引入真偏映射的概念，它使局部紧空间与交换 C*-代数之间的 Gelfand 对应能够以函子化的方式表述。

- 在 §4 中，我们增加了第 13 和第 14 小节，它们讨论实或复可赋范代数中的全纯函数演算，以及无单位元的代数中的全纯函数演算。

- 在 §6 中，连续函数演算的叙述已重新改写（第 6 小节）。第 9 至第 11 小节是新增加的；它们讨论一般 C*-代数中正元素的概念、此类代数的近似单位元，以及 C*-代数对闭双边理想的商。

- §7 论及 Banach 空间自同态的最基本的谱性质，它基本上是新写的；第 7、8 两小节部分地取材于原 §6 的某些小节。正是这一新的一节将作为第 IV 章中更深入地研究 Hilbert 空间自同态的基础。

关于第 II 章：

- 关于 Fourier 变换的 §1 已作了细致的修订，例如讨论 Pontryagin 对偶的函子性的第 6 小节。在新的第 9 小节中，我们把对群 \(\mathbf{R}^n\) 与 \((\mathbf{R}/\mathbf{Z})^n\) 而言最重要的那些命题陈述得更为明确。\(\mathbf{R}^n\) 上的 Fourier 理论，其最自然的框架在于缓增分布理论，这里仅通过其本质性质加以陈述，而不给出证明；这些证明将在第 IV 章中给出。

- 我们最后增加了一张 Fourier 理论一览表，其中汇集了这一理论所涉及的主要函数空间的定义和性质，以及基本的公式。

习题中有许多是新增加的。我们特别指出第 II 章 §1 中习题数量的显著增加，这反映了 Fourier 理论的数学重要性。

